\section{总结与延伸}

\subsection{本章小结}

\begin{importantbox}{四个核心 takeaway}
\begin{enumerate}
    \item 在 LM 里，RL 的 state/action/reward 定义很自然，但 reward 往往非常稀疏。
    \item Naive policy gradient 的形式简单，但方差很大。
    \item Baseline 和 advantage 的本质是降方差，不是改目标。
    \item GRPO 利用了同一 prompt 下多个 response 的 group 结构，所以在 LM 场景里特别合适。
\end{enumerate}
\end{importantbox}


\subsection{最后的提醒}

这一讲其实在传达一个很朴素但重要的判断：
\begin{quote}
如果你能度量某个东西，就可以把它塞进 optimization loop 里，但真正困难的是让这个度量足够稳定、足够密、足够有用。
\end{quote}

这也是为什么 policy gradient 在语言模型里既强大又麻烦：它给了你直接优化目标的入口，也把训练稳定性问题完整暴露了出来。

\subsection{参考}

\begin{itemize}
    \item Schulman et al., Proximal Policy Optimization Algorithms
    \item Deep RL textbook / CS224R policy gradient notes
    \item GRPO / verifiable rewards 相关课程材料
\end{itemize}

\end{document}
